\section{Weighted arctangent families in classical trilogarithms}
\label{xid:wat:sec}
The bilinear integral underlying Section~\ref{cleo:sec:square} reduces to the
product of its half-angle parameters. Inserting the weight $1/\sin x$
changes that structure, but it still leaves an exact evaluation in
classical dilogarithms and trilogarithms. Hyperbolic sums and differences
of the parameters replace their product. This provides a new arithmetic
route at real algebraic parameters, alongside the unit-circle route of
the inverse-hyperbolic companion below. The incoming continuation
\cite{xid:report} supplies these two families; we give their complete
proofs and strengthen its coefficient obstruction to arbitrary
one-variable functions.
\begin{equation}\label{xid:wat:eq:K}
 K(p,q)=\int_0^{\pi/2}
 \frac{\arctan(p\sin x)\arctan(q\sin x)}{\sin x}\,dx,
 \qquad p,q\in\mathbb R.
\end{equation}
The integral in \eqref{xid:wat:eq:K} is ordinary and convergent. Its integrand
is $pq\,x+O(x^3)$ at zero. The principal real arctangent is used
throughout. The polylogarithms below use convergent real series or
absolutely convergent unit-circle boundary series, with the classical
conventions of DLMF~\S25.12~\cite{rw:DLMFPolylog}.

\subsection{The hyperbolic addition formula}

Write
\[
 \chi_j(z)=\frac{\operatorname{Li}_j(z)-\operatorname{Li}_j(-z)}2
          =\sum_{k=0}^{\infty}\frac{z^{2k+1}}{(2k+1)^j}
 \qquad (|z|<1)
\]
for the Legendre chi function. The values used in the next theorem
are real and lie in its ordinary convergent series domain, including
the endpoints $\chi_2(1)=\pi^2/8$ and
$\chi_3(1)=7\zeta(3)/8$.

\begin{theorem}[A real two-term weighted evaluation]
\label{xid:wat:thm:main}
Define the even function
\begin{equation}\label{xid:wat:eq:Fclosed}
 \mathfrak F(s)=\frac{\pi^2|s|}{8}-\frac74\zeta(3)
    +|s|\chi_2(e^{-|s|})+2\chi_3(e^{-|s|}),
 \qquad s\in\mathbb R.
\end{equation}
Its value at zero is $\mathfrak F(0)=0$. The apparent lack of
smoothness from the absolute values is removable: $\mathfrak F$ is
real analytic on $\mathbb R$, and it is equivalently characterized by
\begin{equation}\label{xid:wat:eq:Fode}
 \mathfrak F''(s)=\frac{s}{2\sinh s},\qquad
 \mathfrak F''(0)=\frac12,\qquad
 \mathfrak F(0)=\mathfrak F'(0)=0.
\end{equation}
For every $p,q\in\mathbb R$, put
$\alpha=\operatorname{arcsinh}p$ and
$\beta=\operatorname{arcsinh}q$. Then
\begin{equation}\label{xid:wat:eq:main}
 \boxed{\displaystyle
 K(p,q)=\mathfrak F(\alpha+\beta)-\mathfrak F(\alpha-\beta).}
\end{equation}
In particular the formula includes both diagonals $p=\pm q$,
either zero parameter, and arbitrary signs without branch changes.
\end{theorem}

\begin{proof}
First, for $s>0$ use
$d\chi_j(e^{-s})/ds=-\chi_{j-1}(e^{-s})$ to obtain
\[
 \mathfrak F'(s)=\frac{\pi^2}{8}
              -\chi_2(e^{-s})-s\chi_1(e^{-s}),
 \qquad
 \mathfrak F''(s)=s\chi_0(e^{-s})
                =\frac{s}{2\sinh s}.
\]
The limits at zero follow from the endpoint values of $\chi_2$ and
$\chi_3$ and from
$s\chi_1(e^{-s})=\tfrac{s}{2}\log\coth(s/2)\to0$.
The unique twice-normalized primitive of the real analytic even
function $s/(2\sinh s)$ is real analytic and even. Hence
\eqref{xid:wat:eq:Fclosed} agrees with this primitive on the whole real
line; this proves all assertions in \eqref{xid:wat:eq:Fode}.

Differentiate \eqref{xid:wat:eq:K} once in each parameter. On compact
real parameter sets the differentiated integrand is bounded by an
integrable function, and therefore
\begin{align*}
 K_{pq}(p,q)
 &=\int_0^{\pi/2}
       \frac{\sin x\,dx}
       {(1+p^2\sin^2x)(1+q^2\sin^2x)}\\
 &=\frac{p\operatorname{arcsinh}p/\sqrt{1+p^2}
       -q\operatorname{arcsinh}q/\sqrt{1+q^2}}{p^2-q^2}
       \qquad(p^2\ne q^2).
\end{align*}
For the second equality use partial fractions and
\[
 \int_0^{\pi/2}\frac{\sin x}{1+p^2\sin^2x}\,dx
       =\frac{\operatorname{arcsinh}p}{p\sqrt{1+p^2}},
\]
whose zero-parameter value is $1$ and which follows by
$v=\cos x$.

Let $\widetilde K(\alpha,\beta)=K(\sinh\alpha,\sinh\beta)$.
Multiplication by $\cosh\alpha\cosh\beta$ and the addition formulas
for $\sinh$ now give
\begin{align}\label{xid:wat:eq:hyperbolicmixed}
 \widetilde K_{\alpha\beta}
 &=\frac{\alpha\sinh\alpha\cosh\beta
        -\beta\sinh\beta\cosh\alpha}
       {\sinh^2\alpha-\sinh^2\beta}\notag\\
 &=\frac12\left\{
      \frac{\alpha+\beta}{\sinh(\alpha+\beta)}
     +\frac{\alpha-\beta}{\sinh(\alpha-\beta)}\right\}.
\end{align}
Both quotients on the last line have their removable value $1$ at
zero. Thus the equality holds at $\alpha=\pm\beta$ by continuity.
By \eqref{xid:wat:eq:Fode}, this is the mixed derivative of
$\mathfrak F(\alpha+\beta)-\mathfrak F(\alpha-\beta)$.
Both this expression and $\widetilde K$ vanish on $\alpha=0$ and
$\beta=0$. Integration of their equal mixed derivatives over the
oriented rectangle from $(0,0)$ to $(\alpha,\beta)$ proves
\eqref{xid:wat:eq:main} for every real $\alpha,\beta$.
\end{proof}

For algebraic real $p,q$, the two arguments
$e^{-|\alpha+\beta|}$ and $e^{-|\alpha-\beta|}$ are algebraic.
Indeed $e^\alpha=p+\sqrt{1+p^2}$ and
$e^\beta=q+\sqrt{1+q^2}$ are positive algebraic numbers, and taking
the absolute value in either exponent selects an algebraic number
or its reciprocal. Thus \eqref{xid:wat:eq:main} expresses every such
integral by classical polylogarithms at real algebraic arguments and
logarithms of positive algebraic numbers. It asserts no independence
or minimality of the resulting constants.

The diagonal specializes particularly simply:
\begin{equation}\label{xid:wat:eq:Kdiagonal}
 K(p,p)=\mathfrak F(2\operatorname{arcsinh}p).
\end{equation}
For $\ell=\log(1+\sqrt2)$ and $r=3-2\sqrt2$, this gives the exact
identity
\begin{equation}\label{xid:wat:eq:K11}
 \boxed{\displaystyle
 \int_0^{\pi/2}\frac{\arctan^2(\sin x)}{\sin x}\,dx
 =\frac{\pi^2\ell}{4}-\frac74\zeta(3)
       +2\ell\chi_2(r)+2\chi_3(r).}
\end{equation}
Its numerical value is
$0.718064319741103407036328710693\ldots$; the decimal only
identifies the normalization, while the theorem proves the identity.

An exact weighted counterpart of the manuscript's Cleo integral
also follows. Set $\alpha=\log(1+\sqrt2)$ and
$\beta=\log\bigl((1+\sqrt5)/2\bigr)$. Then
\begin{align}\label{xid:wat:eq:weightedCleo}
 &\int_0^{\pi/2}\frac1{\sin x}
    \arctan^2\!\left(\frac{6\sin x}{3+\cos2x}\right)dx\notag\\
 &\qquad=\mathfrak F(2\alpha)
          +2\mathfrak F(\alpha+\beta)
          -2\mathfrak F(\alpha-\beta)
          +\mathfrak F(2\beta).
\end{align}
To verify the reduction put $s=\sin x\in[0,1]$. Since
$\arctan s+\arctan(s/2)<\pi/2$, the tangent addition formula gives
\[
 \arctan\frac{6\sin x}{3+\cos2x}
       =\arctan(\sin x)+\arctan(\tfrac12\sin x).
\]
Squaring and integrating yields
$K(1,1)+2K(1,1/2)+K(1/2,1/2)$, proving
\eqref{xid:wat:eq:weightedCleo} by \eqref{xid:wat:eq:main}.

\subsection{The weighted inverse-hyperbolic companion}

The same parameter method has a trigonometric form. It gives a
second exact family, including convergent improper endpoint values.

\begin{theorem}[An inverse-hyperbolic bilinear evaluation]
\label{xid:wat:thm:atanh}
For $|r|,|s|\leq1$ define
\[
 L(r,s)=\int_0^{\pi/2}
    \frac{\operatorname{arctanh}(r\sin x)
          \operatorname{arctanh}(s\sin x)}{\sin x}\,dx,
\]
using real inverse hyperbolic tangents and interpreting an endpoint
with $|r|=1$ or $|s|=1$ as an improper integral. Put
\begin{equation}\label{xid:wat:eq:Gclosed}
 \mathfrak G(t)=\frac74\zeta(3)
     -t\operatorname{Im}\chi_2(e^{it})
     -2\operatorname{Re}\chi_3(e^{it}),\qquad |t|\leq\pi.
\end{equation}
The chi values on the unit circle are defined by their absolutely
convergent series. Then $\mathfrak G$ is even, continuous on
$[-\pi,\pi]$, and real analytic on $(-\pi,\pi)$, with
\begin{equation}\label{xid:wat:eq:Gode}
 \mathfrak G''(t)=\frac{t}{2\sin t},\qquad
 \mathfrak G''(0)=\frac12,\qquad
 \mathfrak G(0)=\mathfrak G'(0)=0.
\end{equation}
For $a=\arcsin r$, $b=\arcsin s$, including either sign of either
parameter,
\begin{equation}\label{xid:wat:eq:atanhmain}
 \boxed{\displaystyle
 L(r,s)=\mathfrak G(a+b)-\mathfrak G(a-b).}
\end{equation}
In particular,
\begin{equation}\label{xid:wat:eq:atanh11}
 \boxed{\displaystyle
 \int_0^{\pi/2}\frac{\operatorname{arctanh}^2(\sin x)}{\sin x}\,dx
       =\frac72\zeta(3).}
\end{equation}
\end{theorem}

\begin{proof}
Conjugation of the chi series proves evenness and endpoint
continuity of \eqref{xid:wat:eq:Gclosed}. For $0<t<\pi$,
polylogarithm differentiation gives
\[
 \mathfrak G'(t)=\operatorname{Im}\chi_2(e^{it})
                    -t\operatorname{Re}\chi_1(e^{it}),
 \qquad
 \mathfrak G''(t)=t\operatorname{Im}\chi_0(e^{it})
                 =\frac{t}{2\sin t}.
\]
Here $\operatorname{Re}\chi_1(e^{it})=\tfrac12\log\cot(t/2)$,
so the limit of the first derivative at zero is zero. Also
$\mathfrak G(0)=0$ by $\chi_3(1)=7\zeta(3)/8$.
The normalized even primitive of $t/(2\sin t)$ is real analytic
on $(-\pi,\pi)$, proving \eqref{xid:wat:eq:Gode}.

For $|r|,|s|<1$, differentiation and partial fractions give
\[
 L_{rs}(r,s)=
 \frac{r\arcsin r/\sqrt{1-r^2}
       -s\arcsin s/\sqrt{1-s^2}}{r^2-s^2}
 \qquad(r^2\ne s^2).
\]
The underlying elementary integral is
$\int_0^{\pi/2}\sin x/(1-r^2\sin^2x)\,dx
 =\arcsin r/(r\sqrt{1-r^2})$, again obtained by $v=\cos x$.
For $\widetilde L(a,b)=L(\sin a,\sin b)$ the sine addition formulas
therefore yield
\[
 \widetilde L_{ab}(a,b)
   =\frac12\left\{
       \frac{a+b}{\sin(a+b)}+\frac{a-b}{\sin(a-b)}\right\},
 \qquad |a|,|b|<\pi/2.
\]
The zero quotients take their removable value $1$. This is the
mixed derivative of the right side of \eqref{xid:wat:eq:atanhmain},
and both sides vanish on the two axes. Rectangle integration proves
the identity in the open parameter square.

For all $|r|,|s|\leq1$ the absolute value of the integrand is bounded
by $\operatorname{arctanh}^2(\sin x)/\sin x$. This majorant is
$O(x)$ at zero and $O(\log^2(\pi/2-x))$ at the other endpoint,
so it is integrable. Dominated convergence extends the equality
to the closed square. Finally
$\chi_3(-1)=-7\zeta(3)/8$ and
$\operatorname{Im}\chi_2(-1)=0$ give
$\mathfrak G(\pi)=7\zeta(3)/2$, proving
\eqref{xid:wat:eq:atanh11}.
The endpoint value also follows directly from
$u=\operatorname{arctanh}(\sin x)$:
\[
 \int_0^\infty\frac{u^2}{\sinh u}\,du
    =4\sum_{k=0}^\infty\frac1{(2k+1)^3}
    =\frac72\zeta(3).
\]
\end{proof}

For algebraic $r,s\in[-1,1]$, the chi arguments in this theorem are
algebraic points on the unit circle, since
$e^{ia}=\sqrt{1-r^2}+ir$ and
$e^{ib}=\sqrt{1-s^2}+is$. The theorem has been proved directly over
real parameters; it does not depend on an unstated continuation of
the arctangent formula.

\subsection{An exact obstruction to the one-product ansatz}

The distinction from the unweighted factorization has a short exact proof. Set
\[
 u(p)=\frac{p}{1+\sqrt{1+p^2}},\qquad u(0)=0.
\]
The unweighted integral is a function of $u(p)u(q)$. The weighted
integral has no such factorization into any one-variable function on a
real neighborhood of zero; no regularity assumption on that function is needed.

\begin{proposition}\label{xid:wat:prop:nonfactor}
There is no function $G$ defined on a real neighborhood of zero for which
$K(p,q)=G(u(p)u(q))$ holds for all sufficiently small real $p,q$.
\end{proposition}

\begin{proof}
Suppose such a function exists. Fix a sufficiently small $q_0\ne0$ and
put $c=u(q_0)\ne0$. The local inverse of $u$ is
$p=2u/(1-u^2)$. The asserted identity therefore forces
\[
 G(t)=K\!\left(\frac{2(t/c)}{1-(t/c)^2},q_0\right)
\]
for every sufficiently small real $t$. The integral $K$ is real analytic
near the origin: on compact subdisks of $|p|,|q|<1$, the product of the
arctangent series divided by $\sin x$ converges normally and is dominated
by $C\sin x$ for a parameter-independent constant $C$. Integration of
this normally convergent series proves analyticity despite the apparent
endpoint singularity. Hence the displayed expression forces $G$ to
be real analytic near zero.

For $|p|,|q|<1$, integration of the absolutely convergent arctangent
series gives
\begin{equation}\label{xid:wat:eq:doubleTaylor}
 K(p,q)=\sum_{j,k\geq0}
 \frac{(-1)^{j+k}p^{2j+1}q^{2k+1}}{(2j+1)(2k+1)}
 \frac{4^{j+k}}{(2j+2k+1)\binom{2j+2k}{j+k}}.
\end{equation}
In particular,
\[
 K(p,q)=pq-\frac29(p^3q+pq^3)
       +\text{terms of total degree at least six},
\]
whereas
\[
 u(p)u(q)=\frac{pq}{4}-\frac{p^3q+pq^3}{16}
       +\text{terms of total degree at least six}.
\]
The coefficient of $pq$ would force $G'(0)=4$, and hence the
coefficient of $p^3q$ would have to be $-1/4$, contradicting $-2/9$.
Terms of order at least two in $G$ cannot contain $p^3q$.
\end{proof}

The obstruction concerns this functional ansatz and implies no arithmetic
independence. Theorem~\ref{xid:wat:thm:main} gives a short alternative in
hyperbolic sums and differences. Independent exact checks and signed,
diagonal and endpoint quadratures audit the normalization; the analytic
proofs establish the identities. The ordinary primitive and higher Euler
ladder remain separate integration targets in the editorial ledger.
